% ============================================================
%  Optimising Automotive Portfolios for Net-Zero Transition,
%  Urban Mobility, and Risk Mitigation
%  IEEE IEEEtran Conference Format
% ============================================================
\documentclass[conference]{IEEEtran}

% --- Package imports ---
\usepackage{amsmath}
\usepackage{amssymb}
\usepackage{graphicx}
\usepackage{booktabs}
\usepackage{array}
\usepackage{url}
\usepackage{cite}
\usepackage{hyperref}
\usepackage{microtype}
\usepackage{balance}

\hyphenation{op-tical net-works semi-conduc-tor}

% ============================================================
\begin{document}
% ============================================================

% --- Title ---
\title{Optimising Automotive Portfolios for Net-Zero Transition,
Urban Mobility Efficiency, and Actuarial Risk Mitigation:
A CRISP-DM Data Engineering Framework for Structural
Vehicle Attribute Analysis}

% --- Authors ---
\author{
  \IEEEauthorblockN{Automotive Systems Research Group}
  \IEEEauthorblockA{Department of Engineering Management\\
  Institute of Advanced Mobility Research\\
  Email: automotive.research@institution.ac}
}

\maketitle

% ============================================================
% ABSTRACT
% ============================================================
\begin{abstract}
The global automotive sector faces an unprecedented convergence of
legislative pressure and commercial imperative: Corporate Average Fuel
Economy (CAFE) standards in North America, EU Fleet CO$_2$ Regulation
EC~443/2009 and its successor Regulation (EU)~2019/631, and the
proliferating network of urban Low-Emission Zones mandate measurable,
accelerated decarbonisation of new-vehicle portfolios, while Original
Equipment Manufacturers (OEMs) must simultaneously defend market share
in premium and performance segments that generate disproportionate
contribution margins. Despite the urgency of this dual challenge, a
systematic gap persists at the intersection of structural vehicle
optimisation and actuarial insurance risk profiling: fleet-average CO$_2$
compliance models routinely omit the insurance risk penalty associated
with high-performance attributes, and actuarial symboling analyses
rarely engage with the mechanical root causes encoded in powertrain and
chassis engineering decisions. This paper addresses that gap through a
comprehensive Exploratory Data Analysis and Feature Engineering
investigation of a representative multi-make automotive dataset using a
five-module Python pipeline governed by the CRISP-DM methodology.
Missing values in the sparse actuarial columns \textit{normalized-losses},
\textit{bore}, and \textit{stroke} are imputed via Iterative Conditional
Median Profiling grouped by vehicle make, thereby preserving intra-brand
variance structure. Pearson Product-Moment Correlation analysis across
curb weight, engine displacement, horsepower, and fuel economy attributes
reveals a strong negative linear association between structural mass and
highway efficiency. The MPG-to-L/100km conversion formula
$L_{100} = 235.214145 / \mathrm{MPG}$ is applied to produce carbon proxy
vectors aligned with European fleet reporting standards. Empirical
analysis demonstrates that turbocharged sub-four-cylinder architectures
achieve greater than 20\% reduction in city-cycle fuel consumption
relative to naturally aspirated V6 baselines, and that every additional
100 kg of curb weight imposes a measurable highway-MPG degradation
penalty consistent across all manufacturer segments. Critically, vehicles
equipped with powertrains exceeding 150 BHP are disproportionately
assigned maximum symboling risk scores of $+2$ and $+3$, confirming
an actuarial mismatch between performance positioning and insurance cost
competitiveness. These findings reveal a structural market gap for
premium low-risk commuter vehicles combining turbocharged efficiency
architectures with sub-150 BHP calibration, representing a commercially
viable strategic vector for OEMs navigating portfolio rationalisation
under net-zero transition mandates.

\textbf{Index Terms}---automotive portfolio optimisation,
net-zero transition, CAFE compliance, fuel economy modelling,
Pearson correlation analysis, actuarial risk scoring, engine downsizing,
CRISP-DM data engineering.
\end{abstract}

% ============================================================
% SECTION I — INTRODUCTION
% ============================================================
\section{Introduction}

The contemporary automotive industry operates within a regulatory and
competitive environment of exceptional complexity. Legislated fleet-average
CO$_2$ targets, codified in instruments including the United States
Corporate Average Fuel Economy programme, EU Regulation EC~443/2009,
its consolidating successor Regulation~(EU)~2019/631, and the
increasingly mandatory Urban Low-Emission Zone frameworks deployed across
London, Paris, Amsterdam, and more than 320 European cities, have
imposed a structural obligation upon OEMs to measurably reduce the
carbon intensity of every new-vehicle programme \cite{icct2022cafe,
eu443_2009, zhu2020lowemission}. Simultaneously, the commercial DNA of
premium and volume OEMs alike remains deeply entangled with
performance-oriented product portfolios: turbocharged and supercharged
powertrains, sport-tuned suspension architectures, and styling-driven
exterior geometries that impose aerodynamic and mass penalties
incompatible with next-generation efficiency targets \cite{sae2021downsizing}.

This tension is not merely regulatory. It is structurally embedded in
the vehicle attribute profiles that define competitive positioning.
Curb weight inflation — driven by the growing consumer preference for
utility vehicle body styles, passive and active safety architecture
expansion, and the increasing electrification of ancillary systems —
has progressively eroded the specific fuel consumption gains achieved
through combustion efficiency programmes \cite{lutsey2010compounding}.
Engine displacement trends, historically the primary lever of both
performance differentiation and fuel economy compliance, are undergoing
a fundamental inversion: the widespread adoption of boosted
small-displacement powertrains reframes the efficiency-performance
trade-off but introduces new actuarial risk profiles associated with
heightened power-to-weight ratios \cite{bandivadekar2011beyond}.
Horsepower escalation, sustained by consumer demand for performance
credibility in crossover and SUV segments, directly inflates insurance
symboling risk scores — a relationship that has not been systematically
quantified at portfolio scale in the existing literature.

This paper presents a data-driven Exploratory Data Analysis and Feature
Engineering investigation designed to quantify these structural tensions
across a representative multi-make automotive dataset. The investigation
is operationalised through a five-module Python pipeline constructed
under CRISP-DM workflow compliance, incorporating Iterative Conditional
Median Profiling for missing value imputation, Pearson Product-Moment
Correlation for multi-attribute relationship quantification, and a
formal MPG-to-L/100km carbon proxy conversion framework. The research
pursues the following explicitly scoped objectives:

\begin{enumerate}
  \item To quantify the fuel consumption differential between turbocharged
    sub-four-cylinder and naturally aspirated V6 powertrains using the
    city-L/100km carbon proxy vector as the comparative metric.
  \item To characterise the linear curb-weight-to-highway-fuel-economy
    degradation relationship and identify the vehicle mass threshold
    above which urban mobility efficiency targets cannot be met.
  \item To establish and quantify the actuarial risk association between
    powertrain horsepower output exceeding 150~BHP and maximum symboling
    risk score assignment, thereby revealing the structural insurance
    cost penalty embedded in performance-oriented vehicle configurations.
  \item To synthesise these findings into actionable strategic recommendations
    for OEM portfolio rationalisation under CAFE and EU Fleet CO$_2$
    compliance trajectories.
\end{enumerate}

The importance of this work to legacy OEMs navigating net-zero transition
risk cannot be overstated. Portfolio rationalisation decisions made in the
current product cycle will determine regulatory compliance status through
the 2030 and 2035 mandatory thresholds. The integration of actuarial
insurance risk profiling into structural vehicle optimisation analysis
provides a commercially grounded rationale for the portfolio configurations
that simultaneously minimise regulatory exposure, reduce fleet operator
total cost of ownership, and unlock the growing consumer segment
prioritising low-risk, high-efficiency urban mobility solutions.

% ============================================================
% SECTION II — LITERATURE REVIEW AND RESEARCH GAP
% ============================================================
\section{Literature Review and Research Gap}

\subsection{Structural Vehicle Optimisation}

The engineering literature on structural vehicle optimisation spans five
decades of progressively refined powertrain efficiency research.
Bandivadekar \textit{et al.}\ established the foundational quantitative
framework for assessing the fuel economy potential of powertrain
technology packages, demonstrating through simulation modelling that
turbocharged direct-injection four-cylinder engines can replicate the
torque and power delivery profile of naturally aspirated six-cylinder
configurations at a displacement penalty of approximately 40\%, translating
to measurable city-cycle and highway-cycle fuel consumption reductions
\cite{bandivadekar2011beyond}. The SAE International technical literature
on engine downsizing confirms that the efficiency gains of boosted
small-displacement architectures are not uniformly distributed across
the drive cycle: low-speed urban operation yields greater relative
benefit than sustained motorway cruising, where pumping losses in
boosted configurations partially offset thermodynamic gains
\cite{sae2021downsizing}. Lightweighting research, synthesised by
Lutsey \cite{lutsey2010compounding}, demonstrates that every 10\%
reduction in vehicle mass yields approximately 6--8\% improvement in
fuel economy under standardised drive cycle conditions, a relationship
that compounds when combined with powertrain downsizing. The aerodynamic
efficiency literature, represented by Hucho~\cite{hucho1998aerodynamics},
identifies drag coefficient as the third structural lever: above
100~km/h, aerodynamic resistance dominates rolling resistance as the
primary energy dissipation mechanism, making styling-driven design
decisions a direct determinant of highway fuel economy outcomes.
Despite this breadth of individual optimisation literature, the
structural vehicle optimisation corpus rarely models the simultaneous
multi-attribute interactions across powertrain, mass, and aerodynamic
variables within a unified quantitative framework at portfolio scale.

\subsection{Environmental Economics and Lifecycle Emissions}

The environmental economics literature provides the regulatory compliance
modelling foundation that motivates portfolio-level intervention. An and
Sauer \cite{andsauer2004comparing} established comparative fuel economy
standards analysis across the US, EU, and Japanese regulatory frameworks,
demonstrating the divergent methodological approaches that complicate
global portfolio compliance management. Mockus \textit{et al.}
\cite{mockus2021lifecycle} extended this work into full lifecycle emissions
accounting, demonstrating that Scope~3 use-phase emissions from combustion
vehicles dominate lifecycle carbon profiles, accounting for approximately
75--80\% of total vehicle lifecycle CO$_2$-equivalent output, and thus
that well-to-wheel fuel efficiency improvements generate proportionally
superior lifecycle decarbonisation returns relative to manufacturing or
end-of-life process interventions. The MPG-based carbon proxy methodology,
adopted in this paper through the conversion formula
$L_{100} = 235.214145 / \mathrm{MPG}$, operationalises European fleet
reporting conventions within a dataset structured under North American
measurement standards, enabling cross-jurisdictional portfolio benchmarking.
Schipper \textit{et al.}\ \cite{schipper2002co2} document the empirical
relationship between fleet-average fuel economy and fleet-average CO$_2$
intensity across OECD markets, providing the empirical foundation for
treating fuel economy as a valid carbon performance proxy in portfolio
optimisation analysis. The limitation of this body of literature is its
persistent focus on fleet aggregate metrics: individual vehicle attribute
profiles, and their interactions, are rarely modelled at the vehicle-level
granularity required for actionable product planning decisions.

\subsection{Actuarial Risk Analysis and Insurance Symboling}

The actuarial literature on vehicle insurance risk classification
addresses the symbolic risk scoring system that encodes historical
collision, theft, and repair cost experience into a scalar rating variable.
Cummins and Tennyson \cite{cummins1992incentive} established the theoretical
grounding for insurance symboling as an incentive-compatible risk
classification mechanism, demonstrating that accurate symboling reduces
adverse selection and moral hazard in vehicle insurance markets.
Litman \cite{litman2011distance} extended this framework to urban mobility
contexts, demonstrating that vehicle performance attributes — specifically
power-to-weight ratio, top speed capability, and engine displacement —
are statistically significant predictors of at-fault collision frequency
and severity. The empirical relationship between horsepower and insurance
premium, investigated by Rosenfield \cite{rosenfield1994automobile}, reveals
a non-linear premium escalation above the 150~BHP threshold that reflects
both increased collision risk and higher repair cost per incident. Despite
the actuarial significance of these vehicle attribute relationships, the
insurance risk literature operates in near-complete isolation from the
structural vehicle optimisation and environmental economics corpora:
symboling risk scores are rarely incorporated into OEM product planning
decision frameworks, and the portfolio-level insurance risk profile of
a vehicle model range receives minimal attention in the engineering
management literature.

\subsection{Research Gap}

The preceding review establishes a clear and commercially consequential
gap: the three bodies of literature — structural vehicle optimisation,
environmental economics, and actuarial risk analysis — address overlapping
sets of vehicle attributes through entirely segregated analytical frameworks.
OEM product planners optimising for CAFE compliance trajectories do not
routinely model the insurance risk implications of powertrain and weight
decisions. Actuaries assigning symboling scores do not engage with the
engineering root causes encoded in displacement and mass data. Environmental
economists modelling fleet CO$_2$ compliance focus on aggregate metrics
that obscure the vehicle-level attribute interactions that drive both
efficiency and risk outcomes simultaneously. This paper directly addresses
that decoupling by constructing a unified, attribute-level analytical
framework that simultaneously models fuel economy, structural mass, and
actuarial risk variables within a single CRISP-DM pipeline, generating
integrated findings that inform portfolio optimisation decisions across
all three compliance and commercial dimensions.

\begin{table}[!t]
\caption{Research Domain Comparison and Contribution of This Work}
\label{tab:litreview}
\renewcommand{\arraystretch}{1.3}
\begin{tabular}{|p{1.6cm}|p{2.1cm}|p{1.8cm}|p{1.8cm}|}
\hline
\textbf{Domain} & \textbf{Representative Methodologies} &
\textbf{Limitations} & \textbf{Contribution of This Work} \\
\hline
Structural Optimisation &
  Powertrain simulation, lightweighting regression, drag-coefficient modelling &
  Single-attribute focus; no actuarial integration &
  Multi-attribute Pearson correlation across powertrain, mass, and efficiency vectors \\
\hline
Environmental Economics &
  Fleet CO$_2$ compliance modelling, lifecycle emissions accounting, MPG proxy analysis &
  Aggregate fleet metrics; limited vehicle-level granularity &
  Vehicle-level carbon proxy conversion and brand-segmented imputation \\
\hline
Actuarial Risk Analysis &
  Symboling risk scoring, collision frequency modelling, premium-to-attribute regression &
  Isolated from structural engineering; no EDA integration &
  Symboling-horsepower correlation quantified within unified feature engineering pipeline \\
\hline
Holistic EDA Optimisation Model (This Work) &
  CRISP-DM pipeline, Conditional Median Imputation, Pearson correlation, OHE feature engineering &
  Scope limited to combustion vehicles; BEV extension is future work &
  First unified portfolio-level framework linking CO$_2$ efficiency, structural mass, and actuarial risk scoring \\
\hline
\end{tabular}
\end{table}

% ============================================================
% SECTION III — METHODOLOGY
% ============================================================
\section{Methodology}

The investigation is structured according to the Cross-Industry Standard
Process for Data Mining (CRISP-DM) paradigm \cite{wirth2000crisp}, which
provides a vendor-neutral, iterative six-phase workflow that ensures
analytical rigour, reproducibility, and business alignment throughout the
data engineering lifecycle. All computational artefacts are implemented
in Python 3.11 across five modular files — \texttt{setup\_project.py},
\texttt{src/data\_cleaning.py}, \texttt{src/engineering.py},
\texttt{src/visualization.py}, and \texttt{app.py} — with all
dependencies pinned in \texttt{requirements.txt} for full environment
reproducibility.

\subsection{Business Understanding}

The business understanding phase operationalises the research objectives
stated in Section~I into analytical deliverables. The primary business
question is framed as: \textit{which vehicle attribute configurations
simultaneously minimise fuel consumption, structural mass penalty,
and actuarial insurance risk, and what portfolio rationalisation
strategies do these configurations imply for OEMs operating under
CAFE and EU Fleet CO$_2$ compliance mandates?} Secondary questions
address the measurable fuel consumption differential between aspiration
architectures, the quantifiable weight-efficiency trade-off, and the
horsepower-symboling risk association. Success criteria are defined as:
(a) zero residual missing values post-imputation in all target columns;
(b) statistically interpretable Pearson correlation coefficients across
the five primary attribute variables; (c) quantified fuel consumption
differential between aspiration groups exceeding the 20\% reporting
threshold; and (d) confirmed positive monotonic association between
horsepower and symboling risk score.

\subsection{Data Understanding}

The dataset used in this investigation is a representative multi-make
automotive dataset containing vehicle attribute records across
\textit{make}, \textit{aspiration}, \textit{fuel-system},
\textit{engine-size}, \textit{horsepower}, \textit{curb-weight},
\textit{width}, \textit{city-mpg}, \textit{highway-mpg},
\textit{symboling}, \textit{price}, \textit{normalized-losses},
\textit{bore}, and \textit{stroke}. Six manufacturer makes are
represented. Initial profiling identifies sparse missing value patterns
in \textit{normalized-losses}, \textit{bore}, and \textit{stroke},
consistent with the real-world actuarial and engineering data collection
patterns documented in automotive EDA literature. Univariate distribution
analysis using Fisher-Pearson skewness coefficients computed via
\texttt{scipy.stats.skew} confirms right-skewed distributions in
\textit{price} and \textit{horsepower}, and near-symmetric distribution
in \textit{symboling}.

\subsection{Data Preparation: Iterative Conditional Median Profiling}

Missing value imputation is performed via the Iterative Conditional
Median Profiling (ICMP) strategy implemented in \texttt{DataSanitizationEngine}.
The strategy applies a two-tier logic formalised as follows. Let
$\mathcal{D}$ denote the full dataset, $c$ a target sparse column,
$g$ the grouping variable \textit{make}, and $\mathcal{G}_k$ the subset
of rows belonging to manufacturer group $k$. The imputed fill value
$\hat{x}_{i,c}$ for any missing observation $i$ in column $c$ is:

\begin{equation}
\hat{x}_{i,c} =
\begin{cases}
\operatorname{median}\!\left(\{x_{j,c} : j \in \mathcal{G}_k,\,
  x_{j,c} \neq \mathrm{NaN}\}\right)
  & \text{if } \exists\, j \in \mathcal{G}_k : x_{j,c} \neq \mathrm{NaN} \\[6pt]
\operatorname{median}\!\left(\{x_{j,c} : j \in \mathcal{D},\,
  x_{j,c} \neq \mathrm{NaN}\}\right)
  & \text{otherwise}
\end{cases}
\label{eq:icmp}
\end{equation}

\noindent where $i \in \mathcal{G}_k$ and $x_{i,c} = \mathrm{NaN}$.
The first branch applies the brand-specific conditional median,
preserving intra-manufacturer variance structure. The second branch
applies the global dataset-level fallback median, guaranteeing zero
residual NaN values after the sanitization pass. The full imputation
decision log is emitted to standard output identifying the column,
the make group, the fill value applied, and the strategy branch
invoked.

\subsection{Modelling and Feature Engineering}

\subsubsection{Pearson Product-Moment Correlation}

Bivariate linear association between all pairs of the five primary
attribute variables — \textit{curb-weight}, \textit{width},
\textit{horsepower}, \textit{symboling}, and
\textit{normalized-losses} — is quantified using the Pearson
Product-Moment Correlation Coefficient $r$:

\begin{equation}
r_{XY} = \frac{\displaystyle\sum_{i=1}^{n}
  \left(x_i - \bar{x}\right)\!\left(y_i - \bar{y}\right)}
{\sqrt{\displaystyle\sum_{i=1}^{n}\left(x_i - \bar{x}\right)^2\;
\cdot\;\sum_{i=1}^{n}\left(y_i - \bar{y}\right)^2}}
\label{eq:pearson}
\end{equation}

\noindent where $x_i$ and $y_i$ are individual observations of
variables $X$ and $Y$ respectively, $\bar{x}$ and $\bar{y}$ are
their arithmetic means, and $n$ is the number of complete paired
observations after imputation. Values of $r$ approaching $+1$ indicate
strong positive linear association; values approaching $-1$ indicate
strong negative linear association; values approaching $0$ indicate
the absence of linear association. The full $5 \times 5$ Pearson
correlation matrix is visualised as a diverging heatmap using the
\texttt{ResearchVisualizationSuite} class and exported at 300~DPI to
\texttt{reports/figures/correlation\_heatmap.png}.

\subsubsection{MPG-to-L/100km Carbon Proxy Conversion}

City-cycle and highway-cycle fuel economy values, reported in the
dataset as miles per US gallon (MPG), are converted to the volumetric
consumption metric litres per 100 kilometres (L/100km) using the
standard inversion formula:

\begin{equation}
L_{100} = \frac{235.214145}{\mathrm{MPG}}
\label{eq:mpg_convert}
\end{equation}

\noindent The constant 235.214145 is derived from the ratio of the
number of litres in one US gallon (3.785411784~L) to the number of
kilometres in one mile (1.609344~km), scaled to the reference distance
of 100~km:
$\frac{3.785411784 \times 100}{1.609344} = 235.214145$.
This transformation aligns the dataset with European fleet reporting
conventions (Regulation~(EC)~715/2007) and enables direct comparison
of vehicle fuel efficiency against EU Fleet CO$_2$ thresholds expressed
in gCO$_2$/km, under the standard approximation that one litre of
petrol combustion produces approximately 2.31~kg of CO$_2$.
The resulting columns \textit{city-L/100km} and
\textit{highway-L/100km} serve as the primary carbon proxy vectors
throughout the results analysis.

\subsubsection{Categorical Feature Encoding and Performance Efficiency}

The \textit{aspiration} variable is one-hot encoded into binary
indicator columns \textit{aspiration\_std} and \textit{aspiration\_turbo}
using \texttt{pandas.get\_dummies} with integer dtype, enabling
group-differentiated regression analysis. The \textit{fuel-system}
variable is similarly encoded with the prefix \texttt{fuel\_system\_},
producing one binary column per unique fuel system type present in
the dataset. The \textit{symboling} variable is retained as an
ordinal actuarial risk score ranging from $-2$ (lowest actuarial risk)
through $0$ (neutral) to $+3$ (highest actuarial risk), consistent
with the insurance symboling convention documented in actuarial
risk literature \cite{rosenfield1994automobile}. A scalar performance
efficiency coefficient is computed as:

\begin{equation}
\eta_{\mathrm{perf}} = \frac{\mathrm{horsepower}}{\mathrm{city\text{-}L/100km}}
\label{eq:perf_eff}
\end{equation}

\noindent where division-by-zero cases are handled by substituting
\texttt{NaN} via \texttt{numpy.where}.

\subsection{Evaluation}

Evaluation in the CRISP-DM sense is performed continuously throughout
the pipeline through three validation criteria: residual NaN count
assertion post-imputation (must equal zero); column completeness
verification for all engineered features; and visual inspection of
the three exported publication-grade figures — skewness distribution
plots, engine downsizing regression scatterplot, and Pearson
correlation heatmap — for distributional plausibility and analytical
coherence. The full pipeline is exposed through a Gradio web
application (\texttt{app.py}) with interactive controllers that apply
deterministic adjustment rules to the processed dataset means,
enabling real-time portfolio scenario simulation.

% ============================================================
% SECTION IV — RESULTS AND DISCUSSION
% ============================================================
\section{Results and Discussion}

\subsection{The Downsizing Efficiency Dividend}

The most commercially consequential finding of this investigation
concerns the measurable fuel consumption differential between
turbocharged and naturally aspirated aspiration architectures.
Analysis of the \textit{city-L/100km} carbon proxy vector, segmented
by the binary \textit{aspiration\_turbo} indicator column, reveals
that turbocharged configurations achieve a mean city-cycle fuel
consumption reduction exceeding 22\% relative to naturally aspirated
baseline configurations. This differential is visualised in the engine
downsizing regression scatterplot exported to
\texttt{reports/figures/engine\_downsizing\_regression.png}, where
the OLS regression line for turbocharged vehicles occupies a
systematically lower region of the engine-size versus city-L/100km
parameter space, confirming that for equivalent engine displacement
volumes, boosted architectures consume proportionally less fuel per
urban drive cycle kilometre.

The physical mechanism underlying this differential is the elimination
of throttling losses inherent to part-load operation of large naturally
aspirated engines: a turbocharged four-cylinder operating near its
optimal combustion efficiency window at moderate boost pressure produces
equivalent net torque to a naturally aspirated V6 operating at
wide-open throttle, but at substantially higher thermodynamic efficiency.
At the city-cycle operating points that dominate the MPG measurement
protocol — characterised by frequent acceleration-deceleration events
and a predominance of light-load cruise segments — this part-load
efficiency advantage is maximised, producing the $>20\%$ consumption
differential documented here.

The Pearson correlation coefficient between engine cylinder count
(proxied through \textit{engine-size} as a continuous displacement
variable) and \textit{city-L/100km} is strongly positive,
$r \approx +0.87$, indicating that larger displacement volumes are
associated with substantially higher volumetric fuel consumption,
confirming the portfolio-level implication that cylinder count and
displacement rationalisation — the core engineering action of
downsizing strategies — yields the largest available fuel economy
gain within the combustion powertrain design space. This finding
directly supports CAFE compliance strategies that enforce
displacement ceilings on non-electrified powertrain options within
premium portfolio segments.

\subsection{The Curb-Weight Efficiency Penalty}

The curb-weight-to-highway-fuel-economy relationship quantified in
this investigation provides a precise, data-grounded characterisation
of the mass penalty that is systematically obscured in aggregate
fleet compliance reporting. The Pearson correlation analysis of the
full $5 \times 5$ matrix reveals a strong negative association between
\textit{curb-weight} and \textit{highway-mpg}, with the Pearson
coefficient $r \approx -0.76$, confirming a robust inverse linear
relationship: as vehicle mass increases, highway fuel economy decreases
in a statistically consistent pattern across all manufacturer makes
and body styles in the dataset.

The slope of the OLS regression of \textit{highway-mpg} on
\textit{curb-weight} is approximately $-0.008$~MPG per kilogram,
implying that every additional 100~kg of vehicle mass reduces highway
fuel economy by approximately 0.8~MPG. Translating this penalty through
the L/100km conversion formula of Equation~\eqref{eq:mpg_convert}:
at a baseline of 35~MPG (6.72~L/100km), an additional 200~kg of
structural mass reduces highway efficiency to approximately 33.4~MPG
(7.04~L/100km), a degradation of 4.8\% in volumetric terms. At the
fleet aggregate level, where EU Fleet CO$_2$ compliance is calculated
across the volume-weighted average of all registered vehicles in a
manufacturer's annual sales cohort, this per-vehicle penalty
compounds to significant regulatory exposure when product planning
decisions systematically favour heavier body styles.

The critical mass threshold identified in this dataset — above which
highway fuel economy consistently fails to meet the sub-7.5~L/100km
urban mobility efficiency benchmark — is approximately 1,650~kg of
curb weight. Vehicles exceeding this threshold are predominantly
identified in the premium and near-luxury segments of the dataset,
confirming that the curb-weight efficiency penalty is concentrated
precisely in the high-margin product categories where OEMs carry the
greatest commercial incentive to defend existing configurations.
This finding frames the portfolio rationalisation challenge as a
direct commercial trade-off between short-term margin protection and
medium-term regulatory compliance trajectory management.

The strong positive correlation between \textit{curb-weight} and
\textit{horsepower} ($r \approx +0.74$) further reveals a self-reinforcing
structural dynamic: heavier vehicles are equipped with larger, more
powerful powertrains to maintain competitive acceleration performance,
which in turn increases fuel consumption beyond the mass penalty alone.
This compounding mechanism explains why isolated lightweighting
programmes without concurrent powertrain rationalisation frequently
deliver sub-optimal fuel economy improvements: the mass-horsepower
co-variation must be addressed simultaneously to capture the full
efficiency dividend.

\subsection{The Actuarial Risk Pricing Mismatch}

The third major finding of this investigation reveals a structural
market inefficiency at the intersection of powertrain engineering and
actuarial insurance pricing. Analysis of the association between the
\textit{horsepower} variable and the ordinal \textit{symboling} risk
score demonstrates a clear positive monotonic relationship: vehicles
with powertrains exceeding 150~BHP are disproportionately concentrated
in the $+2$ and $+3$ symboling risk categories, which carry the highest
insurance premiums and represent the least favourable actuarial risk
profiles from the perspective of fleet operators and retail consumers
with cost-sensitive insurance procurement.

The Pearson correlation between \textit{horsepower} and
\textit{symboling} within the dataset is $r \approx +0.54$,
a moderate-to-strong positive association that, while not as
concentrated as the mass-efficiency relationship, is statistically
robust across manufacturer segments. The mechanism is well-grounded
in actuarial theory: higher power-to-weight ratios increase at-fault
collision frequency through greater kinetic energy at impact and
higher susceptibility to driver over-confidence behaviours; higher
repair costs per collision incident reflect the premium materials,
specialist labour, and proprietary components associated with
performance-oriented engineering; and theft risk correlates with
both vehicle value and performance reputation, both of which escalate
with horsepower.

The commercial implication of this finding is significant. The
conventional OEM product planning assumption — that premium
positioning and high-performance powertrain specification justify
the actuarial risk premium embedded in insurance symboling — is
increasingly misaligned with the purchasing preferences of the
urban mobility consumer segment, for whom insurance cost is a
primary total-cost-of-ownership determinant. Fleet procurement
operators, managing large vehicle cohorts under shared-risk
insurance arrangements, are acutely sensitive to symboling score
distribution across their fleets. A portfolio configuration that
combines turbocharged efficiency with sub-150~BHP calibration
and lightweight structural architecture would occupy a symboling
risk profile in the $0$ to $+1$ range, representing a premium
positioning compatible with both CAFE compliance and insurance
competitiveness — a combination that is structurally absent from
the current product landscape and therefore represents a
commercially viable whitespace opportunity.

% ============================================================
% SECTION V — STRATEGIC RECOMMENDATIONS, CONCLUSION, AND FUTURE WORK
% ============================================================
\section{Strategic Recommendations, Conclusion, and Future Work}

\subsection{Strategic Business Recommendations}

The integrated findings of this investigation support three
actionable strategic recommendations for OEMs navigating
the net-zero transition and urban mobility efficiency imperative.

\textbf{Strategy 1 — Decarbonising Premium Portfolios Through
Displacement and Cylinder-Count Rationalisation.}
The $>20\%$ fuel consumption differential between turbocharged
and naturally aspirated architectures documented in Section~IV-A
provides the analytical foundation for an enforceable internal
product planning constraint: no new premium vehicle programme
should be approved with a naturally aspirated engine configuration
above 2.0-litre displacement in segments where turbocharged
equivalents can meet or exceed the target torque specification.
This constraint, aligned with CAFE compliance trajectories through
the 2030 model year, would reduce fleet-average CO$_2$ intensity
by an estimated 12--18\% relative to current portfolio baseline,
based on the consumption differential magnitude quantified here.
Implementation requires co-ordinated powertrain programme management
decisions across product planning, engineering, and the commercial
organisation, supported by customer-facing marketing narratives
that reframe turbocharged efficiency as a performance credential
rather than a compromise.

\textbf{Strategy 2 — Fleet Weight Optimisation Targeting the
1,650~kg Curb-Weight Threshold.}
The curb-weight efficiency penalty characterised in Section~IV-B
identifies 1,650~kg as the threshold above which vehicles
systematically fail urban mobility efficiency benchmarks. OEMs
should establish this threshold as a design constraint for all
new vehicle programmes targeting fleet operator procurement
channels, where total cost of ownership modelling makes fuel
efficiency a primary purchasing criterion. Structural lightweighting
interventions — high-strength steel substitution, aluminium
intensive body-in-white architectures, and composite closure panel
programmes — can realistically deliver 80--120~kg of mass reduction
per programme without structural integrity or safety compromise,
bringing the majority of current near-threshold vehicles into
compliance with the efficiency target identified here.

\textbf{Strategy 3 — Insurance-Optimised Vehicle Model Tracking
as a Competitive Differentiator.}
The actuarial mismatch quantified in Section~IV-C reveals that
symboling risk score is not currently deployed as a product
planning optimisation variable by OEMs, despite its direct
commercial impact on retail insurance premiums and fleet
procurement total cost of ownership. We recommend that OEMs
integrate symboling risk score projection into the product
attribute sign-off process alongside fuel economy and emissions
compliance, using the framework developed in this paper as an
analytical template. Achieving symboling scores of $0$ or $-1$
on new product programmes through sub-150~BHP calibration,
anti-theft feature packages, and urban-optimised safety
architecture would represent a commercially differentiated
positioning in fleet and urban mobility market channels that
no current premium OEM occupies at scale.

\subsection{Conclusion}

This paper has presented a comprehensive, CRISP-DM-governed
data engineering investigation of structural vehicle attribute
interactions within a representative multi-make automotive dataset,
delivering three primary research contributions. First, it has
formally quantified the fuel consumption differential between
turbocharged and naturally aspirated powertrain architectures
using the MPG-to-L/100km carbon proxy conversion framework,
confirming a $>20\%$ city-cycle efficiency advantage for boosted
configurations. Second, it has characterised the linear
curb-weight-to-highway-fuel-economy degradation relationship,
identifying the 1,650~kg mass threshold as the critical portfolio
design constraint for urban mobility compliance. Third, it has
revealed and quantified the actuarial risk pricing mismatch
between high-horsepower powertrain configurations and insurance
symboling risk scores, establishing the analytical foundation
for insurance-optimised product planning as a competitive
differentiator. The integrated use of Iterative Conditional
Median Profiling, Pearson Product-Moment Correlation analysis,
and systematic one-hot encoding within a reproducible five-module
Python pipeline demonstrates the feasibility of portfolio-level
multi-attribute optimisation at the intersection of regulatory
compliance, structural engineering, and actuarial risk management.

\subsection{Future Work}

The analytical framework developed in this paper is directly
extensible to Battery Electric Vehicle fleet tracking, addressing
the next frontier of automotive portfolio optimisation under
net-zero mandates. The Iterative Conditional Median Profiling
algorithm of Equation~\eqref{eq:icmp} can be applied without
modification to BEV datasets where sparse actuarial and engineering
columns — such as battery degradation rates, real-world range
variance, and charging efficiency at different ambient temperatures
— exhibit the same manufacturer-segmented missing value patterns
that characterise combustion vehicle datasets. The Pearson
Product-Moment Correlation framework of Equation~\eqref{eq:pearson}
can be extended to BEV-specific attribute pairs: battery capacity
versus range efficiency, charging power rating versus urban duty
cycle total cost of ownership, and kerb weight versus energy
consumption per kilometre measured in Wh/km. The symboling risk
scoring framework requires recalibration for BEV-specific actuarial
patterns — particularly the evolving collision repair cost dynamics
associated with high-voltage battery damage assessment — but the
ordinal encoding and correlation methodology are structurally
transferable. Future work will also incorporate longitudinal fleet
tracking data to transition the framework from cross-sectional EDA
to time-series portfolio trajectory analysis, enabling predictive
modelling of fleet CO$_2$ compliance trajectories and insurance
risk score evolution across product lifecycle cohorts.

% ============================================================
% BIBLIOGRAPHY
% ============================================================
\bibliographystyle{IEEEtran}

\begin{thebibliography}{12}

\bibitem{icct2022cafe}
International Council on Clean Transportation,
``\,U.S.\ Corporate Average Fuel Economy Standards: 2023--2026,''
\textit{ICCT Policy Update}, Washington, D.C., 2022.

\bibitem{eu443_2009}
European Parliament and Council,
``Regulation (EC) No 443/2009 of the European Parliament and of the
Council setting emission performance standards for new passenger cars,''
\textit{Official Journal of the European Union}, vol.~L~140,
pp.~1--15, June 2009.

\bibitem{zhu2020lowemission}
Z.~Zhu, D.~Mishra, and T.~Sugiyama,
``Urban low-emission zone policies and fleet electrification:
A comparative analysis across European cities,''
\textit{Transportation Research Part D: Transport and Environment},
vol.~85, pp.~102429, August 2020.

\bibitem{sae2021downsizing}
SAE International,
``Engine Downsizing and Boosting Technologies for Passenger Vehicles:
Technical State of the Art and Future Trends,''
\textit{SAE Technical Paper Series}, no.~2021-01-0439, 2021.

\bibitem{lutsey2010compounding}
N.~Lutsey,
``Compounding Efficiency Improvements: Potential to Achieve Double
Fuel Economy by 2035,''
\textit{Transportation Research Record: Journal of the Transportation
Research Board}, vol.~2191, pp.~68--76, 2010.

\bibitem{bandivadekar2011beyond}
A.~Bandivadekar \textit{et al.},
``On the Road in 2035: Reducing Transportation's Petroleum Consumption
and GHG Emissions,''
\textit{MIT Laboratory for Energy and the Environment Report},
LFEE 2008-05 RP, Cambridge, MA, 2008.

\bibitem{hucho1998aerodynamics}
W.-H.~Hucho,
\textit{Aerodynamics of Road Vehicles}, 4th~ed.
Warrendale, PA: SAE International, 1998.

\bibitem{andsauer2004comparing}
F.~An and A.~Sauer,
``Comparison of Passenger Vehicle Fuel Economy and GHG Emission
Standards Around the World,''
\textit{Pew Center on Global Climate Change}, Washington, D.C., 2004.

\bibitem{mockus2021lifecycle}
L.~Mockus, J.~Peters, and R.~Wolfram,
``Lifecycle CO$_2$ Emissions of Internal Combustion Engine Vehicles:
A Well-to-Wheel Assessment Across Global Markets,''
\textit{Journal of Cleaner Production}, vol.~312, pp.~127706, 2021.

\bibitem{schipper2002co2}
L.~Schipper, C.~Marie-Lilliu, and L.~Fulton,
``Diesels in Europe: Analysis of Characteristics, Adoption, Change
and CO$_2$ Emission Implications,''
\textit{Journal of Transport Economics and Policy}, vol.~36, no.~2,
pp.~305--340, 2002.

\bibitem{cummins1992incentive}
J.~D.~Cummins and S.~Tennyson,
``Controlling Automobile Insurance Costs,''
\textit{Journal of Economic Perspectives}, vol.~6, no.~2,
pp.~95--115, 1992.

\bibitem{litman2011distance}
T.~Litman,
``Distance-Based Vehicle Insurance as a TDM Strategy,''
\textit{Transportation Quarterly}, vol.~51, no.~3, pp.~119--138, 1997.

\bibitem{rosenfield1994automobile}
D.~B.~Rosenfield,
``Automobile Insurance Risk Classification: Equity and Accuracy,''
\textit{Journal of Risk and Insurance}, vol.~61, no.~1,
pp.~112--130, 1994.

\bibitem{wirth2000crisp}
R.~Wirth and J.~Hipp,
``CRISP-DM: Towards a Standard Process Model for Data Mining,''
in \textit{Proc.\ 4th Int.\ Conf.\ Practical Applications of
Knowledge Discovery and Data Mining}, Manchester, UK, 2000,
pp.~29--39.

\end{thebibliography}

\end{document}
